% page 97 du fac-simile (image p-096 du scan)
% bloc 162.6 x 242.0 mm ; marge gauche 39.4 mm
\pgc{17.780mm}{0.000mm}{
\l{\cel{50}89}
\l{}
\l{\sou{chuvo}. (\sou{zool}.)Mikra korniko quan onu renkontras en la klosh-}
\l{\cel{3}turmi. - L.\cel{1}\sou{corvus monedula}. F.}
\l{ci. (\sou{gram.}) Pluralo di ca.}
\l{\sou{ciano}. (\sou{kemio)}. Gaso senkolora, produktita da la kombino di}
\l{\cel{3}azoto kun karbo, qua kondutas quale korpo simpla en la reakti}
\l{\cel{3}kemiala. -\cel{1}\sou{Simbolo kemiala}\cel{1}: N : C, C : N. - DEF.}
\l{}
\l{\sou{cianato.(kemio)}. Salo od etero di la cianacido.}
\l{}
\l{\sou{cianhidra}. (\sou{kemio)}. Produktita da la kombino di ciano kun oxo.}
\l{\cel{3}-\cel{1}\sou{Simbolo kemiala}\cel{1}: H. CN.\cel{2}-\cel{3}EF.}
\l{}
\l{\sou{ciatika}. (\sou{anat.}) Qua relatas la grosa nervo qua iras de la}
\l{\cel{3}plexuso sakrumala til la muskuli dopa di la kruro, di la}
\l{\cel{3}suro e di la pedo. - EF.}
\l{}
\l{\sou{cibolo}. (\sou{bot.)}\cel{2}Planto legum-gardenala, analoga a la oniono,}
\l{\cel{3}quan onu uzas koquarte kom kondimento. - L.\cel{1}\sou{allium fistu-}}
\l{\cel{3}\sou{losum}. - DEFIRS.}
\l{}
\l{\sou{ciborio}. (\sou{Liturgio katol}.) Vazo kalico-forma, kun kovrilo,}
\l{\cel{3}qua kontenas la hostii konsakrita por la komuniono di la}
\l{\cel{3}Kristo-Ekleziani. - DEFI.}
\l{}
\l{\sou{cicisbeo}.\cel{2}Amoranto di damo, qua esas quaze oficale agnoskata}
\l{\cel{3}e servas elu publike. - DEFIRS.}
\l{}
\l{\sou{cidro}. Drinkajo ek pomo-suko fermentacita. - DEFIRS.}
\l{}
\l{\sou{cielo.}\cel{2}I. Spaco en qua jiras omna astri.- II. Parto de la spaco}
\l{\cel{3}qua aspektas quale vulto mi-sfera di qua la horizonto esas}
\l{\cel{3}la limito. - FIS.}
\l{}
\l{\sou{cienco}. I. Savo exakta e sistemigita pri ula temo o grupo de}
\l{\cel{3}temi, qua relatas la mondo fizikala od un ek la arti}
\l{\cel{3}liberala (gramatiko, logiko, muziko, e c.) - II. Sistemo}
\l{\cel{3}de savaji, en qua fako de fakti preciza koordinesas e redukt-}
\l{\cel{3}esas a legi(+ leyi). - EFIS.}
\l{}
\l{\sou{cifro}. Tipo o signo qua reprezentas nombro (Araba o Romana). -}
\l{\cel{3}DEFIRS.}
\l{}
\l{\sou{cigano}. Membro di tribui vaganta qui venis de Oriento. - DFIS.}
\l{}
\l{\sou{cigno}. (\sou{zool.})\cel{2}Ucelo palmipeda aquala, analoga a ganso, di qua}
\l{\cel{3}la kolo e la beko esas longa, remarkinda pro la gracio e la}
\l{\cel{3}eleganteso di lua konturi e movi. - L.\cel{1}\sou{cygnus}. - EFIS.}
\l{}
\l{\sou{cikado}.- (\sou{zool}.) Insekto hemiptera qua emisas bruisado akra e}
\l{\cel{3}monotona produktita da la interfroto di du membrani situita}
\l{\cel{3}en la abdomino. - L.\cel{1}\sou{cicada}. - DEFIRS.}
\l{}
\l{\sou{cikatro}. Marko poslasita sur la pelo da vunduro, bruluro, pos}
\l{\cel{3}risanesko. - EFIS.}
}
